\documentclass{article}

\usepackage{tabularx}
\usepackage{booktabs}
\usepackage{float}
\usepackage{longtable}
\usepackage{hyperref}
\hypersetup{
    colorlinks=true,
    linkcolor=blue,
    urlcolor=blue,
}

\title{Reflection and Traceability Report on \progname}

\author{\authname}

\date{}

\input{../Comments}
\input{../Common}

\begin{document}

\maketitle

\newpage
\tableofcontents
\newpage

\plt{Reflection is an important component of getting the full benefits from a
learning experience.  Besides the intrinsic benefits of reflection, this
document will be used to help the TAs grade how well your team responded to
feedback.  Therefore, traceability between Revision 0 and Revision 1 is and
important part of the reflection exercise.  In addition, several CEAB (Canadian
Engineering Accreditation Board) Learning Outcomes (LOs) will be assessed based
on your reflections.}

\section{Changes in Response to Feedback}
\subsection{SRS and Hazard Analysis}

The hazard analysis document was improved through the peer review along with TA feedback in several areas for enhancement, including clarifying the roadmap with implementation timelines and rationales for safety requirements (\href{https://github.com/M9Huynh/technically-functional/issues/144}{\#144}), adding justification for future work prioritization (\href{https://github.com/M9Huynh/technically-functional/issues/150}{\#150}), and removing problematic assumptions about user familiarity with smart devices (\href{https://github.com/M9Huynh/technically-functional/issues/136}{\#136}). Inconsistencies between the SRS and Hazard Analysis were resolved (\href{https://github.com/M9Huynh/technically-functional/issues/129}{\#129}), while the system boundaries and components section was expanded with lower-level explanations to remove ambiguity (\href{https://github.com/M9Huynh/technically-functional/issues/135}{\#135}). The FMEA table was further decomposed to represent detailed failure modes rather than general effects (\href{https://github.com/M9Huynh/technically-functional/issues/142}{\#142}), and hazard categorization was clarified for technical hazards that could impact safety (\href{https://github.com/M9Huynh/technically-functional/issues/153}{\#153}). Additional improvements included adding rationales for each critical assumption (\href{https://github.com/M9Huynh/technically-functional/issues/136}{\#136}), addressing inconsistencies between assumptions and failure modes (\href{https://github.com/M9Huynh/technically-functional/issues/143}{\#143}), and noting where mitigation strategies would be expanded (\href{https://github.com/M9Huynh/technically-functional/issues/152}{\#152}). TA feedback further reinforced the need for roadmap rationales, which was addressed alongside similar peer feedback. One comment regarding bibliography completeness (\href{https://github.com/M9Huynh/technically-functional/issues/151}{\#151}) was determined to be invalid as the citation and reference were already properly included.

\begin{longtable}{>{\raggedright\arraybackslash}p{1.8cm} >{\raggedright\arraybackslash}p{0.8cm} >{\raggedright\arraybackslash}p{4.8cm} >{\raggedright\arraybackslash}p{3.8cm} >{\raggedright\arraybackslash}p{2cm}}
\toprule
\textbf{Source} & \textbf{Issue \#} & \textbf{Issue Description} & \textbf{Changes Made} & \textbf{Corresponding Commit} \\
\midrule
\endhead
\bottomrule
\endfoot

Peer Review & \href{https://github.com/M9Huynh/technically-functional/issues/151}{\#151} & Citation "IppolitoWallace1995" appears but bibliography incomplete & No change - comment invalid as bibliography and in-text citation were complete & None \\
\midrule
Peer Review (Group 14) & \href{https://github.com/M9Huynh/technically-functional/issues/144}{\#144} & Unclear description of how/when safety requirements and hazard testing will be implemented & Added section in roadmap outlining rationale for accepting requirements within capstone timeline and rationales for excluded safety requirements & 939b502 \\
\midrule
Peer Review (Group 14) & \href{https://github.com/M9Huynh/technically-functional/issues/150}{\#150} & Safety requirements prioritized as "first" vs "in the future" without clear justification or timeline & Added clear rationale for future work prioritization & 470F490 \\
\midrule
Peer Review (Group 14) & \href{https://github.com/M9Huynh/technically-functional/issues/136}{\#136} & Missing rationale for each assumption; assumption about user familiarity with smart device should be avoided & Removed assumption about user familiarity; added rationales for each assumption & a103812 \\
\midrule
Peer Review (Group 14) & \href{https://github.com/M9Huynh/technically-functional/issues/143}{\#143} & Inconsistency between assumption about user familiarity and "confusing prompts" failure mode & Removed the initial assumption; will update language of "confusing prompts" & None (different commit) \\
\midrule
Peer Review (Group 14) & \href{https://github.com/M9Huynh/technically-functional/issues/129}{\#129} & Inconsistencies between SRS and Hazard Analysis; hazards need prioritization by severity & Updated consistency between SRS and HA & Reference \#410 PR \\
\midrule
Peer Review (Group 14) & \href{https://github.com/M9Huynh/technically-functional/issues/135}{\#135} & System components need lower-level explanation with sample hazards; authentication module lacks detail & Added changes to HA system boundaries, splitting section to remove ambiguity & 4bd898a \\
\midrule
Peer Review (Group 14) & \href{https://github.com/M9Huynh/technically-functional/issues/142}{\#142} & FMEA rows should represent detailed failure modes rather than effects & Included further decomposition of table within FMEA & Different commit \\
\midrule
Peer Review (Group 14) & \href{https://github.com/M9Huynh/technically-functional/issues/153}{\#153} & Technical hazards could also be safety hazards leading to unsafe exercise & Clarified hazard categorization & 4381828 \\
\midrule
Peer Review (Group 14) & \href{https://github.com/M9Huynh/technically-functional/issues/152}{\#152} & Document lacks detailed mitigation strategies and implementation specifics & Mitigation included; will add section about effectiveness & None (different section) \\
\midrule
TA & N/A & For Section 8 (Roadmap), including rationale would be helpful for which safety requirements to implement and why they are in that time frame & Updated the section with timelines and rationale for why each safety requirement was placed in each term & Similar to \href{https://github.com/M9Huynh/technically-functional/issues/144}{\#144} \\
\bottomrule
\caption{Feedback and changes for SRS and Hazard Analysis}
\label{tab:hazard_feedback}
\end{longtable}

\subsection{Design and Design Documentation}

\subsection{VnV Plan and Report}

\section{Extras}

\plt{Summarize the extras that were tackled by this project.  Extras
can include usability testing, code walkthroughs, user documentation, formal
proof, GenderMag personas, Design Thinking, etc.  Extras should have already
been approved by the course instructor as included in your problem statement.}

\section{Design Iteration (LO11 (PrototypeIterate))}

\plt{Explain how you arrived at your final design and implementation.  How did
the design evolve from the first version to the final version?} 

\plt{Don't just say what you changed, say why you changed it.  The needs of the
client should be part of the explanation.  For example, if you made changes in
response to usability testing, explain what the testing found and what changes
it led to.}

\section{Design Decisions (LO12)}

\plt{Reflect and justify your design decisions.  How did limitations,
 assumptions, and constraints influence your decisions?  Discuss each of these
 separately.}

\section{Economic Considerations (LO23)}

\plt{Is there a market for your product? What would be involved in marketing your 
product? What is your estimate of the cost to produce a version that you could 
sell?  What would you charge for your product?  How many units would you have to 
sell to make money? If your product isn't something that would be sold, like an 
open source project, how would you go about attracting users?  How many potential 
users currently exist?}

\section{Reflection on Project Management (LO24)}

\plt{This question focuses on processes and tools used for project management.}

\subsection{How Does Your Project Management Compare to Your Development Plan}

\plt{Did you follow your Development plan, with respect to the team meeting plan, 
team communication plan, team member roles and workflow plan.  Did you use the 
technology you planned on using?}

\subsection{What Went Well?}

\plt{What went well for your project management in terms of processes and 
technology?}

\subsection{What Went Wrong?}

\plt{What went wrong in terms of processes and technology?}

\subsection{What Would you Do Differently Next Time?}

\plt{What will you do differently for your next project?}

\section{Ethics and Equity (LO18)}

\plt{Describe how your team applied the principle of equity and universal design
to ensure equitable treatment of all stakeholders.}

\section{Reflection on Capstone}

\plt{This question focuses on what you learned during the course of the capstone project.}

\subsection{Which Courses Were Relevant}

\plt{Which of the courses you have taken were relevant for the capstone project?}

\subsection{Knowledge/Skills Outside of Courses}

\plt{What skills/knowledge did you need to acquire for your capstone project
that was outside of the courses you took?}

\end{document}